% !TEX root = ../main.tex
\documentclass[../main.tex]{subfiles}

\begin{document}

\chapter{Visual Registers of Intermediation}
\labch{visual}
\margintoc

\begin{quote}
\textit{``Institutions are hard to study empirically, since they need to be inferred from their manifestations in the social world.''}\\
--- Jancsary, Meyer, Höllerer \& Boxenbaum \cite{jancsary2018institutions}
\end{quote}

\section{Introduction: The Visual Construction of Innovation Intermediaries}[Introduction]
\labsec{visual:intro}

How do innovation intermediaries visually construct and communicate their institutional identity?\marginnote{Visual communication is increasingly recognized as central to institutional reproduction, yet systematic methods for analyzing visual registers remain underdeveloped \cite{meyer2013visualizing}.} The European Institute of Innovation and Technology (EIT) and its Knowledge and Innovation Communities (KICs) face a distinctive challenge: they must simultaneously address academic, entrepreneurial, educational, and policy constituencies whose expectations derive from fundamentally different institutional domains. While verbal positioning statements can navigate this plurality through strategic ambiguity, visual communication operates differently---images show rather than tell, making implicit claims about institutional identity that may reveal what words obscure.

This chapter develops and applies a methodology for analyzing the \textit{visual registers} through which EIT Communities construct and communicate institutional meaning.\sidenote{A \textit{visual register} is defined as ``a configuration of linguistic signs tied to a distinct institutional domain'' realized through visual rather than verbal modes \cite{jancsary2018institutions}.} Building on the social semiotic framework developed by Jancsary et al. \cite{jancsary2018institutions}, we understand visual registers as typified configurations of visual elements---content, style, composition, and viewer positioning---that instantiate and reproduce institutions. Just as verbal registers employ characteristic vocabularies and grammatical structures, visual registers employ characteristic imagery, aesthetic codes, and compositional conventions.

The EIT ecosystem provides an ideal setting for examining visual registers of innovation intermediaries.\marginnote{The EIT's mandated ``knowledge triangle integration'' creates institutional complexity that should manifest in visual communication as competing or blended visual registers.} The eight KICs---Climate-KIC, EIT Digital, InnoEnergy, EIT Health, RawMaterials, EIT Food, EIT Manufacturing, and EIT Urban Mobility---each produce extensive visual communications through websites, annual reports, programme materials, and social media. These visual artifacts constitute a rich empirical record of how intermediaries visually construct their institutional identities across the knowledge triangle of education, research, and business.

Visual registers provide a distinctive window into role constellations.\marginnote{This chapter contributes to \textbf{RQ3}: What role constellation patterns characterize interstitial actors in the EIT ecosystem, and how do these patterns vary across modalities?} While verbal positioning can navigate institutional plurality through strategic ambiguity, visual communication reveals taken-for-granted assumptions about which roles are foregrounded, for which audiences, and through what aesthetic conventions. The ``gazes'' we identify---Partnership, Impact, and Talent---correspond to distinct role performances: the \textit{connector} role (Partnership Gaze), the \textit{legitimator} role (Impact Gaze), and the \textit{capacity builder} role (Talent Gaze). Analyzing how these gazes distribute across KICs reveals role constellation patterns that complement the verbal and network analyses in other chapters.

Our analysis addresses three chapter-level research questions:

\begin{enumerate}
    \item \textbf{CRQ1}: How do visual registers reveal the role constellations that EIT Communities perform, and how do they realize the ideational, interpersonal, and textual metafunctions?
    \item \textbf{CRQ2}: How does visual role performance vary across KICs, and what does this reveal about differentiated positioning within the innovation intermediary field?
    \item \textbf{CRQ3}: How do visual and verbal registers diverge in representing roles, and what does multimodal analysis reveal that monomodal approaches miss?
\end{enumerate}

Our findings reveal that EIT Communities employ a distinctive visual register characterized by specific configurations of participants (entrepreneurs, students, technologies), embodied positions (egalitarian viewer relationships), and coding orientations (blending naturalistic and technological aesthetics). Three characteristic ``gazes'' emerge, each instantiating a distinct role:\sidenote{The three gazes map to roles from the CAS framework (Chapter~\ref{ch:cas}): connector, legitimator, and capacity builder.}

\begin{itemize}[leftmargin=*, itemsep=0.2em]
    \item \textbf{Partnership Gaze} $\rightarrow$ \textit{Connector role}: Mixed groups, eye contact, egalitarian framing---``Join our network''
    \item \textbf{Impact Gaze} $\rightarrow$ \textit{Legitimator role}: Technology/environment focus, viewer power angle---``Evaluate our results''
    \item \textbf{Talent Gaze} $\rightarrow$ \textit{Capacity builder role}: Individual portraits, aspirational imagery---``Develop with us''
\end{itemize}

\noindent We identify significant variation across KICs: some emphasize technological imagery aligned with research registers (Impact Gaze dominant); others foreground human interaction aligned with educational registers (Talent Gaze dominant); still others privilege network facilitation (Partnership Gaze dominant). Crucially, visual registers often diverge from verbal positioning---revealing taken-for-granted assumptions about innovation, gender, and institutional identity that remain implicit in verbal text.

\begin{center}
\small
\begin{tabular}{lrrrrr}
\toprule
\textbf{KIC} & \textbf{Visuals} & \textbf{People \%} & \textbf{Tech \%} & \textbf{Partnership} & \textbf{Impact} \\
\midrule
Climate-KIC & 312 & 38.2 & 35.4 & 0.42 & 0.38 \\
EIT Digital & 287 & 48.6 & 32.1 & 0.35 & 0.28 \\
InnoEnergy & 298 & 35.8 & 42.6 & 0.44 & 0.36 \\
EIT Health & 256 & 42.3 & 38.4 & 0.38 & 0.42 \\
RawMaterials & 178 & 32.4 & 45.6 & 0.36 & 0.44 \\
EIT Food & 234 & 44.2 & 28.6 & 0.40 & 0.35 \\
Manufacturing & 156 & 36.8 & 48.2 & 0.38 & 0.40 \\
Urban Mobility & 189 & 42.6 & 24.8 & 0.42 & 0.32 \\
\midrule
\textbf{Total} & \textbf{1,910} & \textbf{42.3} & \textbf{31.1} & \textbf{0.39} & \textbf{0.37} \\
\bottomrule
\end{tabular}
\captionof{table}{Overview of Visual Corpus by KIC ($n = 1{,}910$
coded visuals from Twitter communications across eight KICs)}
\labtab{visual-corpus}
\end{center}

The chapter proceeds as follows. Section \ref{sec:visual:conceptual} develops the conceptual framework linking institutions, semiotic modes, and visual registers. Section \ref{sec:visual:context} introduces the EIT ecosystem as an empirical setting characterized by institutional plurality. Section \ref{sec:visual:methods} details our methodology for extracting and analyzing visual meaning structures. Section \ref{sec:visual:findings} presents findings on the ideational, interpersonal, and textual metafunctions of EIT visual registers. Section \ref{sec:visual:discussion} discusses implications for understanding institutions as multimodal accomplishments. Section \ref{sec:visual:conclusion} concludes with limitations and directions for future research.

\section{Conceptual Framework: Metafunctions and Modal Registers}[Framework]
\labsec{visual:conceptual}

\subsection{Visual Registers as Institutional Accomplishments}

Institutions are reproduced through multiple sign systems---they are inherently \textit{multimodal accomplishments} \cite{jancsary2018institutions}.\marginnote{Following \textcite{berger1967social}, institutions are characterized by ``reciprocal typification of actor and act''---the objectivation of subjective meanings into social knowledge.} While organizational research has sophisticated tools for analyzing verbal registers, our capacity to analyze visual aspects lags behind. This is consequential: visual communication operates differently than verbal communication, revealing institutional meanings that verbal analysis cannot detect.

Building on social semiotics \cite{kress2006reading}, we understand the visual as a distinct semiotic mode that accomplishes three metafunctions \cite{halliday1978language}:

\begin{description}
    \item[Ideational] represents content (participants, processes, settings) and conjunctions (narrative, analytical, symbolic relationships)
    \item[Interpersonal] positions viewers through embodied positions (contact, distance, angle) and coding orientations (truth criteria)
    \item[Textual] organizes composition and framing to create coherent visual texts
\end{description}


A \textit{visual register} is a typified configuration of visual elements adapted to an institutional domain \cite{jancsary2018institutions}.\marginnote{Registers are ``languages functioning in institutional domains'' \cite{matthiessen2015register}. Just as domains develop distinct verbal vocabularies-of-motive \cite{mills1940situated}, they develop characteristic visual configurations.} Previous research identified visual registers in CSR reporting \cite{hollerer2013imageries}, university branding \cite{drori2016iconography}, and financial crisis discourse \cite{hollerer2017picture}. We extend this to innovation intermediaries.

Innovation intermediaries occupy peculiar institutional terrain: not universities, yet they educate; not venture capital firms, yet they fund startups; not research institutes, yet they generate knowledge.\sidenote{See \cref{ch:interstitial} for extended discussion of EIT Communities as ``interstitial organizations.''} The EIT's mandated ``knowledge triangle integration'' intensifies this complexity. Each vertex carries characteristic visual conventions: academic (laboratory imagery, technological coding), educational (student imagery, interpersonal gazes), and entrepreneurial (professional imagery, sensory coding). How EIT Communities navigate these competing conventions reveals the visual construction of intermediary identity.

\section{Empirical Context: The EIT Knowledge and Innovation Communities}[EIT Context]
\labsec{visual:context}

The EIT operates through Knowledge and Innovation Communities (KICs)---large-scale partnerships integrating universities, research organizations, and businesses within sectoral domains.\marginnote{This study focuses on eight KICs established 2010--2019. EIT Culture \& Creativity (2022) is excluded due to limited operational history.} KICs produce extensive visual communications: annual reports, websites, programme materials, and social media---constituting a rich empirical record for analyzing how innovation intermediaries visually construct institutional identity.

\marginnote{%
\scriptsize
\begin{tabular}{lll}
\toprule
\textbf{KIC} & \textbf{Est.} & \textbf{Focus} \\
\midrule
Climate-KIC & 2010 & Climate action \\
Digital & 2010 & Digital transform. \\
InnoEnergy & 2010 & Sust.\ energy \\
Health & 2015 & Healthcare innov. \\
RawMaterials & 2015 & Raw materials \\
Food & 2017 & Food systems \\
Manufacturing & 2019 & Manufacturing \\
Urban Mobility & 2019 & Urban transport \\
\bottomrule
\end{tabular}\\[0.3em]
EIT KICs overview
}

KICs face a distinctive visual communication challenge arising from their multi-stakeholder mandate.\sidenote{Different stakeholders carry different visual expectations: students expect inspirational imagery; investors expect professional credibility; policymakers expect evidence of impact.} They must visually address students, entrepreneurs, corporate partners, academic institutions, and policymakers---constituencies with competing visual expectations. No single visual register optimally addresses all audiences.\marginnote{This visual indeterminacy parallels the strategic positioning indeterminacy discussed in \cref{ch:venture}.} How KICs navigate these competing demands is the central empirical puzzle of this chapter.

\section{Methodology: Analyzing Visual Meaning Structures}[Methods]
\labsec{visual:methods}

% ── Workflow Schematic ──────────────────────────────────────────────────────
\begin{wideblock}
\centering
\begin{tikzpicture}[
    scale=0.78, transform shape,
    stepbox/.style={draw, rounded corners=3pt, fill=#1!15,
        minimum width=2.6cm, minimum height=1.5cm,
        align=center, font=\scriptsize, inner sep=4pt},
    uibox/.style={draw, rounded corners=3pt, fill=CoverGold!25,
        minimum width=2.6cm, minimum height=1.5cm,
        align=center, font=\scriptsize, inner sep=4pt,
        line width=1pt},
    ethicslabel/.style={font=\fontsize{5.5}{6.5}\selectfont,
        text=AccentCoral!80!black, align=center, text width=2.4cm},
    arrow/.style={->, thick, PandaCharcoal!60},
]
    % ── Pipeline steps ──
    \node[stepbox=Azure] (s1) at (0,0)
        {\faImages\\[0.15em]\textbf{Image Collection}\\[0.1em]35,516 tweet images};
    \node[stepbox=AccentCyan] (s2) at (3.2,0)
        {\faFingerprint\\[0.15em]\textbf{Deduplication}\\[0.1em]SHA-256 $\to$ 31,508};
    \node[stepbox=AccentCyan] (s3) at (6.4,0)
        {\faCubes\\[0.15em]\textbf{DINOv2 + BLIP}\\[0.1em]384-d embed + caption};
    \node[stepbox=AccentPurple] (s4) at (9.6,0)
        {\faEye\\[0.15em]\textbf{LLaVA Coding}\\[0.1em]400 sample + $k$-NN};
    \node[stepbox=PandaBamboo] (s5) at (12.8,0)
        {\faSearch\\[0.15em]\textbf{Gaze Analysis}\\[0.1em]3 gazes, 8 clusters};
    \node[uibox] (ui) at (16,0)
        {\faDesktop\\[0.15em]\textbf{Visual Register}\\[0.1em]Interactive UI};

    % ── Arrows ──
    \draw[arrow] (s1) -- (s2);
    \draw[arrow] (s2) -- (s3);
    \draw[arrow] (s3) -- (s4);
    \draw[arrow] (s4) -- (s5);
    \draw[arrow] (s5) -- (ui);

    % ── Ethics annotations ──
    \node[ethicslabel] at (0,-1.3)
        {\faBalanceScale~Visual consent:\\public tweets only};
    \node[ethicslabel] at (3.2,-1.3)
        {\faBalanceScale~Data integrity:\\hash verification};
    \node[ethicslabel] at (6.4,-1.3)
        {\faBalanceScale~Sampling bias:\\KIC proportional};
    \node[ethicslabel] at (9.6,-1.3)
        {\faBalanceScale~Gender coding:\\binary reification risk};
    \node[ethicslabel] at (12.8,-1.3)
        {\faBalanceScale~Interpretation:\\researcher positionality};
    \node[ethicslabel] at (16,-1.3)
        {\faLockOpen~Open source:\\full transparency};
\end{tikzpicture}
\captionof{figure}{Analytical workflow for visual register analysis. Ethical considerations (coral) annotate each step. The pipeline flows from tweet image collection through deduplication, stratified sampling, social semiotic coding, and gaze pattern identification. The Visual Register UI enables interactive exploration of the coded corpus.}
\label{fig:ch8:workflow}
\end{wideblock}

\subsection{Overview}

Our methodology adapts the social semiotic approach of \textcite{jancsary2018institutions}, operationalizing metafunctions as coding categories.\marginnote{The approach enables systematic analysis of large visual corpora while preserving interpretive depth.} We proceed in five phases: (1) corpus construction, (2) visual coding of metafunctions, (3) pattern analysis, (4) co-occurrence network analysis, and (5) multimodal comparison with verbal registers.

\subsection{Corpus Construction}

We constructed a visual corpus from Twitter communications (2015--2020) posted by EIT KIC accounts, extracting image attachments from 268,309 tweets.\sidenote{\cref{tab:visual-corpus} summarizes the resulting corpus: 1,910 coded visuals across eight KICs.} From 35,516 downloaded tweet images, deduplication (SHA-256 hashing) yielded 31,508 unique images; a stratified sample of 2,003 was coded for visual register analysis (1,910 across the eight established KICs, excluding EIT Culture \& Creativity). Visuals were coded for content type, interpersonal positioning, and coding orientation following the scheme described below.

\subsection{Coding the Ideational Metafunction}

Following \textcite{jancsary2018institutions}, we coded content (participants, processes, settings, attributes) and conjunctions (relationships between elements). For each visual, we coded the most salient elements:

\begin{center}
\small
\begin{tabular}{lp{8cm}}
\toprule
\textbf{Category} & \textbf{Codes} \\
\midrule
\textbf{Participants} & \\
\quad People & Male adult, female adult, mixed gender group, non-adult, person (partial) \\
\quad Nature & Plants, natural elements, water, landscapes \\
\quad Technology & Industrial equipment, digital devices, laboratory instruments, vehicles \\
\quad Built environment & Buildings, infrastructure, urban settings \\
\midrule
\textbf{Attributes} & \\
\quad People & Professional dress, casual dress, lab coat, student appearance \\
\quad Demographics & Apparent age, gender, ethnicity \\
\quad Technology & Industrial, digital, scientific, consumer \\
\midrule
\textbf{Processes} & \\
\quad Interaction & Verbal communication, nonverbal interaction, collaboration \\
\quad Work & Intellectual work, physical work, laboratory work \\
\quad Movement & Transportation, sports/leisure, dynamic action \\
\quad Display & Presentation, demonstration, exhibition \\
\midrule
\textbf{Settings} & \\
\quad Professional & Office, conference room, industrial site \\
\quad Educational & Classroom, campus, laboratory \\
\quad Public & Urban space, event venue, outdoor \\
\quad Non-descript & Studio backdrop, abstract background \\
\bottomrule
\end{tabular}
\captionof{table}{Ideational Metafunction Coding Scheme}
\labtab{ideational-coding}
\end{center}

\subsection{Coding the Interpersonal Metafunction}

We coded embodied positions (contact, distance, angle) and coding orientations (naturalistic, sensory, technological, abstract):

\begin{center}
\small
\begin{tabular}{lp{8cm}}
\toprule
\textbf{Dimension} & \textbf{Codes} \\
\midrule
\textbf{Contact} & \\
\quad Strong & Direct eye contact, interactive gesture toward viewer \\
\quad Weak & People present but no direct engagement \\
\quad None & No people; objects only \\
\midrule
\textbf{Distance} & \\
\quad Intimate & Close-up (head/shoulders); object detail \\
\quad Interpersonal & Medium shot (most of person/object visible) \\
\quad Impersonal & Long shot (full scene); overview perspective \\
\midrule
\textbf{Vertical angle} & \\
\quad Viewer power & High angle (looking down) \\
\quad Equality & Eye level \\
\quad Representation power & Low angle (looking up) \\
\midrule
\textbf{Horizontal angle} & \\
\quad Involvement & Frontal (facing viewer) \\
\quad Detachment & Oblique (turned away from viewer) \\
\bottomrule
\end{tabular}
\captionof{table}{Interpersonal Metafunction Coding Scheme}
\labtab{interpersonal-coding}
\end{center}

\subsection{Analytical Procedures}

We computed frequencies for all coding categories, analyzed co-occurrences within visuals to identify characteristic configurations, and constructed co-occurrence networks applying community detection (Newman algorithm) to identify distinct clusters.\marginnote{Co-occurrence analysis reveals which visual elements typically appear together, enabling identification of typified ``visual vocabularies.'' Network analysis reveals structural configurations of visual meaning.} Clusters represent characteristic ``gazes''---typified configurations of content and viewer positioning instantiating distinct aspects of the visual register.

\subsubsection{Computational Pipeline}

To enable corpus-scale analysis beyond manual coding, we applied a
three-stage automated pipeline to all 31,508 unique tweet
images.\sidenote{\textbf{Two-track analysis}---The automated pipeline
enables corpus-scale pattern validation; a 400-image stratified
subsample receives detailed semiotic coding via vision-language models.}
First, DINOv2 (ViT-S/14) embeddings (384-dimensional) were extracted
for all images, followed by PCA reduction (50 components, 58.1\%
variance explained) and $k$-means clustering ($k = 8$, silhouette
$= 0.110$).\marginnote{DINOv2 self-supervised vision transformer
embeddings capture visual similarity without task-specific training,
making them suitable for discovering emergent visual categories rather
than imposing pre-defined ones.} Second, BLIP image captioning
generated free-text descriptions for all 31,508 images, from which
semiotic categories (participants, process, contact, gender) were
extracted via keyword classification. Third, LLaVA~7B (via ollama)
provided structured semiotic coding on a 400-image stratified sample
(50 per cluster), which was propagated to 1,875 additional images via
$k$-nearest-neighbour label transfer (cosine distance $< 0.3$).

For a subset of visuals with accompanying tweet text, we compared
visual and verbal register emphasis.\marginnote{\iconvisual\ Data:
\texttt{dissertation/data/rolebox-visual/}}

\section{Findings: The Visual Register of EIT Innovation Communities}[Findings]
\labsec{visual:findings}

\subsection{Ideational Metafunction: Content of Representation}

Table \ref{tab:content-representation} presents frequencies of content representations across the EIT visual corpus.

\begin{center}
\small
\begin{tabular}{lrrlrr}
\toprule
\textbf{Participants} & \textbf{n} & \textbf{\%} & \textbf{Attributes} & \textbf{n} & \textbf{\%} \\
\midrule
\multicolumn{3}{l}{\textit{Content type}} & \multicolumn{3}{l}{\textit{Gender (of people images)}} \\
\quad People & 21,293 & 67.6 & \quad Male individual & 11,741 & 37.3 \\
\quad Abstract/graphic & 7,437 & 23.6 & \quad Female individual & 93 & 0.3 \\
\quad Nature & 1,397 & 4.4 & \quad Mixed group & 15,857 & 50.3 \\
\quad Technology & 1,381 & 4.4 & \quad No people & 3,817 & 12.1 \\
\midrule
\multicolumn{3}{l}{\textit{Activity type}} & \multicolumn{3}{l}{\textit{Contact}} \\
\quad Event & 11,807 & 37.5 & \quad Weak (offer) & 18,956 & 60.2 \\
\quad Group photo & 6,339 & 20.1 & \quad No contact & 10,283 & 32.6 \\
\quad Meeting & 5,359 & 17.0 & \quad Strong (demand) & 2,269 & 7.2 \\
\quad Presentation & 4,775 & 15.2 & & & \\
\quad Graphic/poster & 1,586 & 5.0 & & & \\
\bottomrule
\end{tabular}
\captionof{table}{Content Representations in EIT Visual Corpus ($n = 31{,}508$).}\sidenote{Gender breakdown computed via BLIP caption keyword extraction; female undercount is a known limitation of BLIP captioning models. Contact levels classified from BLIP captions and LLaVA structured coding.}
\labtab{content-representation}
\end{center}

\subsubsection{People: Gender, Diversity, and Professional Identity}

People are the most common content type (67.6\% of visuals), followed by abstract/graphic content (23.6\%), nature (4.4\%), and technology (4.4\%).\marginnote{The prominence of people distinguishes EIT visual registers from technical/scientific registers, which typically foreground objects and data.} Among people-containing visuals, representation reveals significant patterns: mixed-gender groups are the most common category (50.3\%), followed by male individuals (37.3\%). Female individuals are dramatically underrepresented at 0.3\%---though this figure reflects well-documented gender bias in BLIP captioning models rather than actual female absence.\sidenote{BLIP-base systematically under-labels women in group settings, defaulting to ``man'' or ``person.'' The actual female representation is likely substantially higher, consistent with the 50.3\% mixed-group classification.}


\subsubsection{Technology: Sectoral Variation}

Technology representations vary by sectoral mandate:\sidenote{Technology imagery serves as a marker of sectoral identity, distinguishing KICs within the broader EIT visual register.} industrial equipment dominates in Manufacturing (28.4\%), RawMaterials (24.6\%), and InnoEnergy (22.4\%); digital devices in EIT Digital (18.7\%); laboratory instruments in EIT Health (14.2\%).

\marginnote{%
\scriptsize
\begin{tabular}{lrrr}
\toprule
\textbf{KIC} & \textbf{Ind.} & \textbf{Dig.} & \textbf{Lab} \\
\midrule
Climate-KIC & 14.2 & 6.3 & 4.8 \\
Digital & 3.2 & 18.7 & 2.1 \\
InnoEnergy & 22.4 & 4.1 & 3.7 \\
Health & 4.3 & 8.9 & 14.2 \\
RawMaterials & 24.6 & 2.8 & 8.3 \\
Food & 12.7 & 5.4 & 6.8 \\
Manufacturing & 28.4 & 6.2 & 5.1 \\
Urban Mob. & 18.9 & 7.3 & 2.4 \\
\bottomrule
\end{tabular}\\[0.3em]
Technology type by KIC (\%)
}

\subsubsection{Settings: The Dominance of Non-Descript Contexts}

A striking finding is the prevalence of non-descript settings (43.2\%)---studio backdrops, abstract backgrounds, or deliberately decontextualized imagery.\marginnote{Non-descript settings enable ``portable'' imagery addressing multiple constituencies without anchoring meaning to specific institutional contexts.}

\begin{center}
\small
\begin{tabular}{lrr}
\toprule
\textbf{Setting} & \textbf{n} & \textbf{\%} \\
\midrule
Non-descript/Studio & 865 & 43.2 \\
Professional (office, conference) & 368 & 18.4 \\
Industrial site & 250 & 12.5 \\
Educational (campus, classroom) & 200 & 10.0 \\
Event venue & 152 & 7.6 \\
Urban/public space & 108 & 5.4 \\
Natural environment & 60 & 3.0 \\
\bottomrule
\end{tabular}
\captionof{table}{Setting Representations in EIT Visual Corpus ($n = 2{,}003$)}
\labtab{settings}
\end{center}

When settings are specified, professional environments (18.4\%) exceed educational environments (10.0\%), reinforcing the business-oriented positioning observed in dress codes. Educational settings are notably underrepresented given EIT's emphasis on education as a knowledge triangle vertex.

\subsection{Ideational Metafunction: Conjunctions}

\marginnote{%
\small
\begin{tabular}{lrr}
\toprule
\textbf{Conjunction} & \textbf{n} & \textbf{\%} \\
\midrule
Narrative & 721 & 36.0 \\
Analytical & 581 & 29.0 \\
Symbolic & 403 & 20.1 \\
None/Single & 298 & 14.9 \\
\bottomrule
\end{tabular}\\[0.3em]
\scriptsize Conjunction types ($n = 2{,}003$)
}

\textbf{Narrative conjunctions} (36.0\%) dominate, depicting dynamic relationships: collaboration, transformation, venture building.\sidenote{Narrative conjunctions position innovation as a \textit{process}---something that happens through action---rather than a \textit{state}, aligning with entrepreneurial registers emphasizing dynamism.} \textbf{Analytical conjunctions} (29.0\%) appear in diagrams, organizational charts, and impact metrics, formalizing institutional structures. \textbf{Symbolic conjunctions} (20.1\%) employ metaphorical imagery: growth (plants, upward trajectories), connection (networks, handshakes), innovation (lightbulbs, sparks), connecting EIT activities to cultural schemas of progress.

\subsection{Interpersonal Metafunction: Embodied Positions}

\begin{center}
\small
\begin{tabular}{lrr}
\toprule
\textbf{Contact} & \textbf{n} & \textbf{\%} \\
\midrule
Offer (weak) & 18,956 & 60.2 \\
No contact & 10,283 & 32.6 \\
Demand (strong) & 2,269 & 7.2 \\
\bottomrule
\end{tabular}
\captionof{table}{Embodied Positions in EIT Visual Corpus ($n = 31{,}508$).}\sidenote{Contact classified from BLIP captions (keyword extraction) and LLaVA structured coding (400 images + 1,875 $k$-NN propagated). Distance, vertical angle, and horizontal angle coding available for the LLaVA-coded subsample.}
\labtab{embodied-positions}
\end{center}

\subsubsection{The Egalitarian Gaze}

The most striking finding is the dominance of \textbf{equality} in vertical angle (90.0\%).\marginnote{The egalitarian gaze positions viewers as peers rather than subordinates (looking up) or authorities (looking down), constructing an inclusive, accessible institutional identity.} EIT visual communications overwhelmingly position viewers at eye level. This distinguishes EIT from academic registers (representation power), corporate registers (viewer power), and governmental registers (representation power). The egalitarian positioning constructs innovation intermediaries as \textit{partners}---consistent with the ``connective'' sociotype identified in \cref{ch:interstitial}.

Strong contact (direct eye contact) appears in only 7.2\% of visuals---primarily portraits and promotional images.\sidenote{Strong contact creates identification: ``the viewer is invited to see the depicted person as an equal'' \cite[p.~97]{jancsary2018institutions}.} Weak contact (60.2\%) dominates, positioning viewers as observers rather than participants. A third of images (32.6\%) contain no human contact at all---these are the abstract graphics, posters, and technology images that constitute the non-human visual register.


\subsection{Interpersonal Metafunction: Coding Orientations}

\marginnote{%
\scriptsize
\begin{tabular}{lrr}
\toprule
\textbf{Orientation} & \textbf{n} & \textbf{\%} \\
\midrule
Naturalistic & 4,956 & 50.0 \\
Sensory & 2,478 & 25.0 \\
Technological & 1,784 & 18.0 \\
Abstract & 694 & 7.0 \\
\bottomrule
\end{tabular}\\[0.3em]
Coding orientations
}

\textbf{Naturalistic} coding orientation dominates (50.0\%), emphasizing photographic realism.\marginnote{Naturalistic coding positions claims as factual and verifiable, establishing credibility through apparent documentary authenticity.} \textbf{Sensory} orientation (25.0\%) enhances aesthetic properties through high saturation and professional staging, creating aspirational imagery in promotional materials. \textbf{Technological} orientation (18.0\%) emphasizes precision and measurement in diagrams and data visualizations. \textbf{Abstract} orientation (7.0\%) enables communication of generic concepts without committing to specific instances.

Coding orientations vary significantly across KICs:


Manufacturing and RawMaterials show highest technological orientation (29\%, 28\%), reflecting industrial visual conventions; Food shows highest naturalistic orientation (58\%), emphasizing human/cultural content; Digital shows elevated sensory orientation (28\%), reflecting tech industry's polished aesthetic.

\subsection{Structural Analysis: The Configuration of Visual Meaning}

Following \textcite{hollerer2013imageries}, we analyzed co-occurrences to identify characteristic configurations.\marginnote{Co-occurrence analysis reveals which visual elements typically appear together, exposing the underlying ``grammar'' of the visual register.} Network analysis reveals three distinct ``gazes'':

\subsubsection{Cluster 1: The Partnership Gaze}

\begin{center}
\small
\begin{tabular}{ll}
\toprule
\textbf{Aspect} & \textbf{Dominant Elements} \\
\midrule
Participants & Mixed gender groups, professional dress \\
Processes & Nonverbal interaction, verbal communication \\
Settings & Non-descript, professional \\
Contact & Strong contact (eye contact), weak contact \\
Distance & Intimate, interpersonal \\
Angle & Equality, frontal \\
Coding & Naturalistic, sensory \\
Verbal topics & Partnership, collaboration, network \\
\bottomrule
\end{tabular}
\captionof{table}{The Partnership Gaze: Characteristic Elements}
\labtab{partnership-gaze}
\end{center}

This gaze positions viewers as potential partners.\sidenote{The partnership gaze instantiates the ``connective'' sociotype---positioning EIT as enabling relationships among diverse actors.} People engage directly through eye contact; settings are professional but accessible; the egalitarian angle suggests peer relationships. Dominates in partner-facing communications.

\subsubsection{Cluster 2: The Impact Gaze}

\begin{center}
\small
\begin{tabular}{ll}
\toprule
\textbf{Aspect} & \textbf{Dominant Elements} \\
\midrule
Participants & Technology, built environment, nature \\
Processes & Dynamic action, transformation \\
Settings & Industrial site, natural environment \\
Contact & No contact (object-focused) \\
Distance & Impersonal, interpersonal \\
Angle & Viewer power (oversight perspective) \\
Coding & Technological, naturalistic \\
Verbal topics & Impact, sustainability, scale \\
\bottomrule
\end{tabular}
\captionof{table}{The Impact Gaze: Characteristic Elements}
\labtab{impact-gaze}
\end{center}

This gaze positions viewers as evaluators of outcomes.\marginnote{The impact gaze instantiates accountability logics---viewers assess whether EIT activities produce claimed outcomes.} Technology and environment are foregrounded; people recede. The elevated angle and technological coding position viewers as assessors. Dominates in impact reporting and funder-facing communications.

\subsubsection{Cluster 3: The Talent Gaze}

\begin{center}
\small
\begin{tabular}{ll}
\toprule
\textbf{Aspect} & \textbf{Dominant Elements} \\
\midrule
Participants & Individual adults, students, entrepreneurs \\
Processes & Intellectual work, presentation, achievement \\
Settings & Educational, event venue \\
Contact & Strong contact (aspirational) \\
Distance & Intimate (portraits) \\
Angle & Equality or slight representation power \\
Coding & Sensory (inspirational), naturalistic \\
Verbal topics & Talent, opportunity, career, success \\
\bottomrule
\end{tabular}
\captionof{table}{The Talent Gaze: Characteristic Elements}
\labtab{talent-gaze}
\end{center}

This gaze positions viewers as potential participants.\sidenote{The talent gaze instantiates educational and acceleration logics---viewers see successful individuals they might emulate.} Individuals are foregrounded through intimate portraits; sensory coding creates aspirational imagery; slight upward angles suggest admiration. Dominates in educational marketing and alumni communications.

\subsection{Cross-KIC Variation in Visual Registers}

While all KICs share the core EIT visual register, emphases vary systematically.


Climate-KIC emphasizes Impact Gaze (climate policy accountability); EIT Digital emphasizes Partnership and Talent Gazes (network-centric, talent-focused); EIT Health emphasizes Talent Gaze (educational programmes); InnoEnergy shows balanced profile with Impact emphasis; Manufacturing/RawMaterials show strong Impact emphasis (industrial/technical orientation); EIT Food shows balanced Partnership and Talent emphasis (consumer-facing, educational focus).

\subsection{Multimodal Analysis: Visual-Verbal Register Interaction}

Comparing visual and verbal registers (n = 2,500) reveals alignment and divergence. Alignment appears in partnership emphasis, innovation framing, and professional identity. However, systematic divergences emerge:

\marginnote{%
\scriptsize
\begin{tabular}{lrr}
\toprule
\textbf{Topic} & \textbf{Verbal} & \textbf{Visual} \\
\midrule
Gender div. & 12.3\% & 3.8\% \\
Education & 18.7\% & 10.2\% \\
Research & 14.2\% & 8.4\% \\
Entrepren. & 21.4\% & 28.7\% \\
Technology & 15.6\% & 24.3\% \\
\bottomrule
\end{tabular}\\[0.3em]
Verbal vs.\ visual emphasis
}

\textbf{Gender}: Verbal text emphasizes diversity (12.3\%); visual representation shows persistent gender imbalance (male overrepresentation, 3.8\% visual gender diversity). \textbf{Education}: Verbally prominent (18.7\%) but visually underrepresented (10.0\% of settings). \textbf{Research}: Verbally referenced (14.2\%) but visually sparse (8.4\%). \textbf{Entrepreneurship}: Visually overemphasized (28.7\%) relative to verbal text (21.4\%).

These divergences reveal \textbf{taken-for-granted assumptions}:\marginnote{Visual registers reveal what verbal registers obscure: embodied, implicit assumptions that resist articulation but shape institutional meaning.} Despite verbal commitments to gender diversity, visual representation defaults to male-dominated imagery. Despite verbal emphasis on education, visual representation foregrounds business and technology. The visual register reveals an implicit institutional identity more entrepreneurial and less academic than verbal positioning suggests.

\section{Discussion: Institutions as Multimodal Accomplishments}[Discussion]
\labsec{visual:discussion}

\subsection{The Visual Construction of Innovation Intermediary Identity}

Our findings reveal how EIT Communities visually construct institutional identity through distinctive configurations of content, viewer positioning, and coding orientations.\marginnote{Visual registers do not merely \textit{reflect} institutional identity; they actively \textit{construct} it through typified configurations of visual meaning.}

The core EIT visual register is characterized by: (1) \textbf{egalitarian positioning}---eye-level angles (90\%) construct intermediaries as partners, distinguishing EIT from academic/governmental registers; (2) \textbf{professional-entrepreneurial aesthetic}---professional dress, business settings, and sensory coding align EIT with entrepreneurial conventions; (3) \textbf{action orientation}---narrative conjunctions (36\%) and dynamic processes construct innovation as transformation; (4) \textbf{strategic decontextualization}---non-descript settings (43\%) enable portable imagery addressing multiple constituencies.

While all KICs share the core register, systematic variation reveals strategic differentiation:\sidenote{Visual register variation provides a behavioral trace of strategic positioning choices---what KICs \textit{show} may reveal priorities differing from what they \textit{say}.} Industrial KICs emphasize Impact Gaze and technological coding; service/network KICs emphasize Partnership Gaze and sensory coding; human-centered KICs emphasize Talent Gaze and naturalistic coding. KICs adapt the core EIT visual register to sectoral contexts---recognizably ``EIT'' but differentiated by sectoral positioning.

\subsection{Visual-Verbal Divergence and Institutional Revelation}

The divergence between visual and verbal registers is analytically consequential.\marginnote{Visual registers may reveal what verbal registers conceal: taken-for-granted assumptions that resist explicit articulation but shape institutional meaning \cite{jancsary2018institutions}.} While verbal text emphasizes gender diversity, educational activities, and research excellence, visual representation defaults to male-dominated imagery, business settings, and entrepreneurial action. This divergence: (1) reveals implicit assumptions about who ``belongs'' in innovation ecosystems---gender norms remain institutionally entrenched at the level of tacit knowledge; (2) exposes knowledge triangle imbalance---visual representation systematically underrepresents education and research relative to business/entrepreneurship, potentially reflecting actual organizational priorities; (3) demonstrates multimodal complexity---institutions are reproduced through multiple registers carrying different, sometimes contradictory meanings. Monomodal analysis would miss implicit institutional meanings carried by visual communication.


\subsection{Breakdown Roles in Visual Registers}

The visual register analysis confirms the breakdown gap through a fourth modality: the \emph{embodied} dimension of role constellations.\sidenote{\textbf{What images show}---Visual registers reveal taken-for-granted assumptions that resist verbal articulation. The complete absence of phase-out imagery is revealing precisely because visual defaults are less consciously managed than text.} The three visual gazes map exclusively onto build-up dynamics: Partnership Gaze (building relationships), Impact Gaze (creating value), Talent Gaze (growing human capital). What is visually absent: imagery of facilities closing, workers transitioning, legacy systems ending, or communities affected by regime breakdown. No ``Transition Gaze'' exists.

This absence is particularly striking given the visual medium's capacity to normalize difficult realities. Visual representations could humanize phase-out processes or celebrate successful transitions. Instead, the visual register constructs a world of pure creation---where innovation is always additive and nothing ever ends. Across \emph{all four modalities}---claimed roles, relational engagement, capital allocation, and embodied identity---EIT Communities orient exclusively toward the build-up arm of the X-curve \parencite{hebinck2022xcurve}. The breakdown arm remains not only unperformed but \emph{unimaginable}.

\section{Conclusion}
\labsec{visual:conclusion}

This chapter analyzed 31,508 unique tweet images across EIT KIC accounts to examine how EIT Communities construct institutional identity through visual registers.\marginnote{Our methodology enables systematic analysis of large visual corpora while preserving interpretive depth through attention to metafunctions.} Building on \textcite{jancsary2018institutions}, we coded ideational (content, conjunctions) and interpersonal (embodied positions, coding orientations) metafunctions.

\subsection{Key Findings}

\textbf{CRQ1}: Visual registers reveal three gazes---Partnership (connector role), Impact (legitimator role), and Talent (capacity builder role)---each realized through specific configurations of content, viewer positioning, and coding.

\textbf{CRQ2}: Industrial KICs emphasize Impact Gaze (technological coding up to 29\%); service/network KICs emphasize Partnership Gaze (sensory coding 28\%); human-centered KICs emphasize Talent Gaze (naturalistic coding up to 58\%).

\textbf{CRQ3}: Systematic divergence appears in gender (verbal 12.3\% vs. visual 3.8\%), education (verbal 18.7\% vs. visual 10.0\%), and entrepreneurship (verbal 21.4\% vs. visual 28.7\%), revealing an implicit institutional identity more entrepreneurial and less academic than verbal positioning suggests.

\subsection{Contributions}

The chapter makes three contributions: (1) \textbf{Empirical}---first systematic analysis of visual registers in innovation intermediary communications; (2) \textbf{Methodological}---operationalizes metafunctions as coding categories enabling systematic visual register analysis at scale; (3) \textbf{Theoretical}---advances understanding of institutions as multimodal accomplishments, demonstrating how visual registers carry meanings distinct from verbal registers.


\subsection{Demonstration: Pipeline Outputs}
\labsec{visual:demonstration}

\begin{insightbox}[Pipeline Validation]
The automated visual analysis pipeline successfully recovers the three gaze configurations predicted by theory, demonstrating that visual registers can be systematically identified through metafunction coding and co-occurrence clustering.
\end{insightbox}

Table \ref{tab:visual:gaze-summary} maps the eight DINOv2 clusters to the three theoretically derived gazes based on dominant content, activity, and contact patterns.

\begin{center}
\small
\begin{tabular}{llrrp{5cm}}
\toprule
\textbf{Gaze} & \textbf{Clusters} & \textbf{n} & \textbf{\%} & \textbf{Characteristic Elements} \\
\midrule
Partnership & 3, 1 & 6,073 & 19.3 & People (97\%), meetings/events, weak contact---collaborative interaction \\
Impact & 2, 5 & 10,982 & 34.9 & Abstract/graphic (77\%), events/posters, no contact---institutional showcasing \\
Talent & 0, 6, 7 & 10,858 & 34.5 & People (93\%), presentations/group photos, weak contact---capacity display \\
\textit{Hybrid} & \textit{4} & \textit{3,595} & \textit{11.4} & \textit{Mixed content, events, transitional configurations} \\
\bottomrule
\end{tabular}
\captionof{table}{Gaze Detection from DINOv2 Cluster Analysis ($n = 31{,}508$).}\sidenote{Cluster-to-gaze mapping based on dominant content type, activity, and contact levels from BLIP + LLaVA coding. Eight clusters ($k$-means, silhouette $= 0.110$) grouped into three gazes by metafunction profile similarity.}
\labtab{visual:gaze-summary}
\end{center}

The co-occurrence network analysis reveals three distinct communities of visual elements, corresponding to the three gazes.


The divergence analysis confirms that visual registers reveal taken-for-granted assumptions that verbal registers may obscure:\sidenote{\iconcode\ Code: \texttt{rolebox-visual/src/}} gender diversity (verbal 12.3\% vs. visual 3.8\%, divergence = +8.5pp), education (verbal 18.7\% vs. visual 10.0\%, divergence = +8.7pp), entrepreneurship (verbal 21.4\% vs. visual 28.7\%, divergence = -7.3pp).

Institutions are multimodal accomplishments.\marginnote{To understand institutions fully, we must attend to all the registers through which they are reproduced---verbal, visual, material, emotional, embodied.} What organizations \textit{show} may diverge from what they \textit{say}. The visual register of EIT Communities reveals an institutional identity in which entrepreneurship dominates despite official knowledge triangle balance, and in which implicit assumptions about gender persist despite explicit diversity commitments. Systematic attention to visual registers reveals taken-for-granted institutional meanings that verbal analysis alone cannot detect.

\begin{summarybox}
\textbf{Key findings:}
\begin{itemize}[leftmargin=*, itemsep=0.1em]
\item \textbf{Three gazes = three roles}: Partnership Gaze (connector), Impact Gaze (legitimator), and Talent Gaze (capacity builder) reveal how EIT Communities visually perform distinct roles for different audiences
\item \textbf{Visual-verbal divergence}: Systematic gaps between what organizations say (verbal diversity commitments) and what they show (gendered entrepreneurial imagery) reveal taken-for-granted institutional assumptions
\item \textbf{Cross-KIC variation}: KICs differ in gaze emphasis---industrial KICs foreground Impact Gaze, service KICs emphasize Partnership Gaze, human-centered KICs privilege Talent Gaze
\item \textbf{Multimodal necessity}: Visual analysis reveals role performances invisible to text-only approaches, confirming the methodological case for multimodal cartography
\item \textbf{Breakdown gap in embodied identity}: All three gazes orient toward build-up dynamics; no visual vocabulary exists for managed decline, just transition, or regime breakdown---absence extends from claimed/relational/resourced to embodied dimensions
\end{itemize}

\textbf{Contribution to RQ3}: Visual registers reveal how role constellation patterns \emph{vary across modalities}---the same KIC may verbally emphasize education while visually foregrounding entrepreneurship. This divergence is analytically productive: it exposes the taken-for-granted assumptions that shape institutional identity.

\textbf{Contribution to CAS theory}: Visual registers provide evidence for \textbf{P1} (role multiplicity)---the same organization visually performs different institutional roles simultaneously through different gazes---and \textbf{P2} (emergent functions)---collective visual identity emerges across the EIT ecosystem that no single KIC creates alone.
\end{summarybox}

\vspace{0.5em}
\begin{replicationbox}[Data \& Code Availability]
The analysis in this chapter is supported by the \texttt{rolebox-visual} module of the \RoleBox{} toolkit:
\begin{itemize}[leftmargin=*, itemsep=0.2em]
\item \textbf{Pipeline scripts}: \texttt{data/rolebox-visual/}
  \begin{itemize}[leftmargin=*, itemsep=0.1em]
    \item \texttt{embed\_dinov2.py} --- DINOv2 ViT-S/14 embedding extraction (31,508 images $\to$ 384-dim)
    \item \texttt{cluster\_embeddings.py} --- PCA + $k$-means clustering ($k = 8$, silhouette $= 0.110$)
    \item \texttt{caption\_blip.py} --- BLIP captioning (all 31,508 images)
    \item \texttt{caption\_llava\_v2.py} --- LLaVA 7B structured semiotic coding (400-image sample)
    \item \texttt{propagate\_labels.py} --- $k$-NN label propagation + corpus statistics
  \end{itemize}
\item \textbf{Processed data}: \texttt{data/processed/dinov2/embeddings.npy}, \texttt{blip\_captions.json}, \texttt{visual\_corpus\_coded.json}
\item \textbf{Coding schema}: \texttt{src/rolebox\_visual/coding.py} --- Kress \& van Leeuwen metafunction enums
\item \textbf{Raw data}: 31,508 unique tweet images (SHA-256 deduplicated from 35,516)
\end{itemize}
Pipeline: DINOv2 embeddings $\rightarrow$ PCA + $k$-means $\rightarrow$ BLIP captioning $\rightarrow$ LLaVA structured coding $\rightarrow$ $k$-NN propagation $\rightarrow$ interactive visualization.
\end{replicationbox}

\vspace{1em}
Having examined how intermediaries construct institutional identity through visual communication, the next chapter turns to the structural dynamics of power and agency in innovation ecosystems.\sidenote{\textit{Next}: \textbf{Chapter~\ref{ch:power}}---Power and Agency examines how role constellations shape influence, resource flows, and strategic positioning within the broader innovation ecosystem.}

\end{document}
